% =====================================================================
% INSERT FOR PAPER 0  (Arithmetic Expression Geometry 0)
%
% Source file:  paper-0/sections/04b-pole-covariance.tex
% Integrated 2026-10-20 as Subsection 4.4 of Section 4 ("Ripple Geometry: The
%   Projective Dual"), immediately after Subsection 4.3 ("Ripple histories of
%   right-expanded expressions") and before the status subsection, via
%   % =====================================================================
% INSERT FOR PAPER 0  (Arithmetic Expression Geometry 0)
%
% Source file:  paper-0/sections/04b-pole-covariance.tex
% Integrated 2026-10-20 as Subsection 4.4 of Section 4 ("Ripple Geometry: The
%   Projective Dual"), immediately after Subsection 4.3 ("Ripple histories of
%   right-expanded expressions") and before the status subsection, via
%   % =====================================================================
% INSERT FOR PAPER 0  (Arithmetic Expression Geometry 0)
%
% Source file:  paper-0/sections/04b-pole-covariance.tex
% Integrated 2026-10-20 as Subsection 4.4 of Section 4 ("Ripple Geometry: The
%   Projective Dual"), immediately after Subsection 4.3 ("Ripple histories of
%   right-expanded expressions") and before the status subsection, via
%   % =====================================================================
% INSERT FOR PAPER 0  (Arithmetic Expression Geometry 0)
%
% Source file:  paper-0/sections/04b-pole-covariance.tex
% Integrated 2026-10-20 as Subsection 4.4 of Section 4 ("Ripple Geometry: The
%   Projective Dual"), immediately after Subsection 4.3 ("Ripple histories of
%   right-expanded expressions") and before the status subsection, via
%   \input{sections/04b-pole-covariance}  at the end of
%   paper-0/sections/04-ripple-geometry.tex.
%
% Labels required from the host paper (all present):
%   eq:p0-mobius-map, def:p0-ripple-pencil, thm:p0-ripple-pencil-formula,
%   prop:p0-path-ripple-duality, eq:p0-projective-factorization,
%   eq:p0-bruhat, sec:p0-projective-unification.
% New labels introduced here (checked unique in the host paper):
%   subsec:p0-pole-covariance, lem:p0-pole-covariance,
%   prop:p0-pole-rigidity, rem:p0-pole-bruhat, rem:p0-pole-forward.
%
% Scope note.  Paper 0's scope file permits "the dual reading (poles,
% horocycle pullback, ripple pencils), reciprocals, infinity, fixed points,
% and the matrix synthesis", and excludes "projective bivaluations, quotient
% towers, condensation".  Everything below is therefore phrased in poles and
% in the Bruhat decomposition only; the quantitative comparison of two
% pencils tangent at different poles is left to Paper IV by a prospective
% reference.
%
% Claim status: Lemma and Proposition proved; the two Remarks record
% consequences and a prospective interface, not new theorems.
%
% Numbers when inserted in the current build of Paper 0 (Section 4 ends at
% statement 4.5): Lemma 4.6, Proposition 4.7, Remarks 4.8 and 4.9, and the
% subsection is 4.4, shifting the status subsection to 4.5.
% =====================================================================

\subsection{The pole of a prefix: covariance and rigidity}
\label{subsec:p0-pole-covariance}

Section~\ref{sec:p0-ripples} introduced the pole $p_F=F^{-1}(\infty)$ of a
single projective operator, and the ripple history of a word
$F_1,\dots,F_n$ carries the pole sequence $p_k=G_k^{-1}(\infty)$ of the prefixes
$G_k=F_k\circ\cdots\circ F_1$.  It was observed there that
\emph{the common tangency point of the pencil may move with $k$}.  The two
elementary statements below make that observation exact: they record how a
pole transforms under composition, and they identify the histories along which
it does not move.  Neither statement is new mathematics; they are recorded
because the mobility of the tangency point is otherwise left as a picture.

Throughout, a M\"obius map is as in \eqref{eq:p0-mobius-map}, and
$B_\infty=\Stab_{\PGL_2(K)}(\infty)$ is the affine Borel sector of
Section~\ref{sec:p0-projective-unification}.

\begin{lemma}[Covariance of the pole]
\label{lem:p0-pole-covariance}
Let $P,Q$ be nondegenerate M\"obius maps and put $p_P:=P^{-1}(\infty)$.  Then
\begin{equation}\label{eq:p0-pole-covariance}
  p_{P\circ Q}=Q^{-1}\bigl(p_P\bigr).
\end{equation}
Consequently
\begin{equation}\label{eq:p0-pole-constancy-criterion}
  p_{P\circ Q}=p_Q
  \quad\Longleftrightarrow\quad
  p_P=\infty
  \quad\Longleftrightarrow\quad
  P\in B_\infty .
\end{equation}
\end{lemma}

\begin{proof}
$(P\circ Q)^{-1}=Q^{-1}\circ P^{-1}$, so
$p_{P\circ Q}=(Q^{-1}\circ P^{-1})(\infty)=Q^{-1}\bigl(p_P\bigr)$, which is
\eqref{eq:p0-pole-covariance}.  For the criterion: since $Q^{-1}$ is injective
on the projective line, $Q^{-1}(p_P)=Q^{-1}(\infty)$ holds exactly when
$p_P=\infty$, and $P^{-1}(\infty)=\infty$ holds exactly when
$P(\infty)=\infty$, that is $C_P=0$, that is $P\in B_\infty$.
\end{proof}

Lemma~\ref{lem:p0-pole-covariance} says that the pole moves by the action of the
\emph{next} operator on the projective line: the new pole is the old one moved by
$Q^{-1}$.  It also says that a step leaves the pole unchanged exactly when the
step is affine.  Applied to a history, the two statements give the rigidity
criterion.

\begin{proposition}[Rigidity of the pole along a history]
\label{prop:p0-pole-rigidity}
Let $F_1,\dots,F_n$ be nondegenerate M\"obius sections in chronological order,
$G_k=F_k\circ\cdots\circ F_1$, and $p_k=G_k^{-1}(\infty)$.  The following are
equivalent:
\begin{enumerate}
  \item $p_k=p_{k-1}$ for $k=1,\dots,n$;
  \item $F_k\in B_\infty$ for $k=1,\dots,n$, i.e. every lower-left entry
  vanishes;
  \item $p_k=\infty$ for $k=1,\dots,n$.
\end{enumerate}
Equivalently, the pole sequence is not constantly $\infty$ exactly when the
history leaves the affine sector, and then the first index at which $C_k\ne0$ is
also the first index at which the tangency point is finite.
\end{proposition}

\begin{proof}
Apply Lemma~\ref{lem:p0-pole-covariance} with $P=F_k$ and $Q=G_{k-1}$: the step
$G_{k-1}\to G_k$ leaves the pole unchanged exactly when $F_k\in B_\infty$.  This
gives (1)$\Leftrightarrow$(2).  Since $G_0=\id$ and $p_0=G_0^{-1}(\infty)
=\infty$, the equivalence (2)$\Leftrightarrow$(3) follows by induction: if
every $F_k$ is affine then every prefix is affine, so $C_k=0$ and
$p_k=\infty$; conversely if some $C_k\ne0$ then $p_k=-D_k/C_k\ne\infty$ by the definition of the pole in
Section~\ref{sec:p0-ripples}.
\end{proof}

\begin{remark}[The pole is read off one factor of the factorization]
\label{rem:p0-pole-bruhat}
When $C_k\ne0$, the factorization \eqref{eq:p0-projective-factorization} of
$G_k$ reads
\[
  G_k
  =T_{A_k/C_k}\circ D_{\Delta_k/C_k^{2}}\circ J\circ T_{D_k/C_k},
\]
and the pole is
$p_k=G_k^{-1}(\infty)=-D_k/C_k$, i.e. the negative of the translation parameter
of the innermost factor.  Thus of the four factors, the one that determines the
tangency point of the pencil is the rightmost, while the value orbit
$G_k(z_0)$ depends on the whole word.  This is the precise sense in which the
pole and the point are independent data of one prefix, and it is the reading of
the cell decomposition \eqref{eq:p0-bruhat} in terms of tangency: the pole is
$\infty$ exactly in the affine cell, and it is finite exactly in the projective
cell.
\end{remark}

\begin{remark}[Prospective interface]
\label{rem:p0-pole-forward}
When a history leaves $B_\infty$, consecutive pencils $\mathcal R_k$ and
$\mathcal R_{k+1}$ are based at different points of $\partial_\infty\mathbb H^2$
and therefore have no common tangency point; the families cannot be nested, and
the pair carries a nonzero relative datum.  Quantifying that datum is not part
of this paper, whose scope fixes it to the poles and pencils themselves.  A
stabilizer-valued transport on the projective quotient that measures it is
constructed in Paper~IV; this remark is a prospective interface only and is not
used anywhere in the present paper.
\end{remark}
  at the end of
%   paper-0/sections/04-ripple-geometry.tex.
%
% Labels required from the host paper (all present):
%   eq:p0-mobius-map, def:p0-ripple-pencil, thm:p0-ripple-pencil-formula,
%   prop:p0-path-ripple-duality, eq:p0-projective-factorization,
%   eq:p0-bruhat, sec:p0-projective-unification.
% New labels introduced here (checked unique in the host paper):
%   subsec:p0-pole-covariance, lem:p0-pole-covariance,
%   prop:p0-pole-rigidity, rem:p0-pole-bruhat, rem:p0-pole-forward.
%
% Scope note.  Paper 0's scope file permits "the dual reading (poles,
% horocycle pullback, ripple pencils), reciprocals, infinity, fixed points,
% and the matrix synthesis", and excludes "projective bivaluations, quotient
% towers, condensation".  Everything below is therefore phrased in poles and
% in the Bruhat decomposition only; the quantitative comparison of two
% pencils tangent at different poles is left to Paper IV by a prospective
% reference.
%
% Claim status: Lemma and Proposition proved; the two Remarks record
% consequences and a prospective interface, not new theorems.
%
% Numbers when inserted in the current build of Paper 0 (Section 4 ends at
% statement 4.5): Lemma 4.6, Proposition 4.7, Remarks 4.8 and 4.9, and the
% subsection is 4.4, shifting the status subsection to 4.5.
% =====================================================================

\subsection{The pole of a prefix: covariance and rigidity}
\label{subsec:p0-pole-covariance}

Section~\ref{sec:p0-ripples} introduced the pole $p_F=F^{-1}(\infty)$ of a
single projective operator, and the ripple history of a word
$F_1,\dots,F_n$ carries the pole sequence $p_k=G_k^{-1}(\infty)$ of the prefixes
$G_k=F_k\circ\cdots\circ F_1$.  It was observed there that
\emph{the common tangency point of the pencil may move with $k$}.  The two
elementary statements below make that observation exact: they record how a
pole transforms under composition, and they identify the histories along which
it does not move.  Neither statement is new mathematics; they are recorded
because the mobility of the tangency point is otherwise left as a picture.

Throughout, a M\"obius map is as in \eqref{eq:p0-mobius-map}, and
$B_\infty=\Stab_{\PGL_2(K)}(\infty)$ is the affine Borel sector of
Section~\ref{sec:p0-projective-unification}.

\begin{lemma}[Covariance of the pole]
\label{lem:p0-pole-covariance}
Let $P,Q$ be nondegenerate M\"obius maps and put $p_P:=P^{-1}(\infty)$.  Then
\begin{equation}\label{eq:p0-pole-covariance}
  p_{P\circ Q}=Q^{-1}\bigl(p_P\bigr).
\end{equation}
Consequently
\begin{equation}\label{eq:p0-pole-constancy-criterion}
  p_{P\circ Q}=p_Q
  \quad\Longleftrightarrow\quad
  p_P=\infty
  \quad\Longleftrightarrow\quad
  P\in B_\infty .
\end{equation}
\end{lemma}

\begin{proof}
$(P\circ Q)^{-1}=Q^{-1}\circ P^{-1}$, so
$p_{P\circ Q}=(Q^{-1}\circ P^{-1})(\infty)=Q^{-1}\bigl(p_P\bigr)$, which is
\eqref{eq:p0-pole-covariance}.  For the criterion: since $Q^{-1}$ is injective
on the projective line, $Q^{-1}(p_P)=Q^{-1}(\infty)$ holds exactly when
$p_P=\infty$, and $P^{-1}(\infty)=\infty$ holds exactly when
$P(\infty)=\infty$, that is $C_P=0$, that is $P\in B_\infty$.
\end{proof}

Lemma~\ref{lem:p0-pole-covariance} says that the pole moves by the action of the
\emph{next} operator on the projective line: the new pole is the old one moved by
$Q^{-1}$.  It also says that a step leaves the pole unchanged exactly when the
step is affine.  Applied to a history, the two statements give the rigidity
criterion.

\begin{proposition}[Rigidity of the pole along a history]
\label{prop:p0-pole-rigidity}
Let $F_1,\dots,F_n$ be nondegenerate M\"obius sections in chronological order,
$G_k=F_k\circ\cdots\circ F_1$, and $p_k=G_k^{-1}(\infty)$.  The following are
equivalent:
\begin{enumerate}
  \item $p_k=p_{k-1}$ for $k=1,\dots,n$;
  \item $F_k\in B_\infty$ for $k=1,\dots,n$, i.e. every lower-left entry
  vanishes;
  \item $p_k=\infty$ for $k=1,\dots,n$.
\end{enumerate}
Equivalently, the pole sequence is not constantly $\infty$ exactly when the
history leaves the affine sector, and then the first index at which $C_k\ne0$ is
also the first index at which the tangency point is finite.
\end{proposition}

\begin{proof}
Apply Lemma~\ref{lem:p0-pole-covariance} with $P=F_k$ and $Q=G_{k-1}$: the step
$G_{k-1}\to G_k$ leaves the pole unchanged exactly when $F_k\in B_\infty$.  This
gives (1)$\Leftrightarrow$(2).  Since $G_0=\id$ and $p_0=G_0^{-1}(\infty)
=\infty$, the equivalence (2)$\Leftrightarrow$(3) follows by induction: if
every $F_k$ is affine then every prefix is affine, so $C_k=0$ and
$p_k=\infty$; conversely if some $C_k\ne0$ then $p_k=-D_k/C_k\ne\infty$ by the definition of the pole in
Section~\ref{sec:p0-ripples}.
\end{proof}

\begin{remark}[The pole is read off one factor of the factorization]
\label{rem:p0-pole-bruhat}
When $C_k\ne0$, the factorization \eqref{eq:p0-projective-factorization} of
$G_k$ reads
\[
  G_k
  =T_{A_k/C_k}\circ D_{\Delta_k/C_k^{2}}\circ J\circ T_{D_k/C_k},
\]
and the pole is
$p_k=G_k^{-1}(\infty)=-D_k/C_k$, i.e. the negative of the translation parameter
of the innermost factor.  Thus of the four factors, the one that determines the
tangency point of the pencil is the rightmost, while the value orbit
$G_k(z_0)$ depends on the whole word.  This is the precise sense in which the
pole and the point are independent data of one prefix, and it is the reading of
the cell decomposition \eqref{eq:p0-bruhat} in terms of tangency: the pole is
$\infty$ exactly in the affine cell, and it is finite exactly in the projective
cell.
\end{remark}

\begin{remark}[Prospective interface]
\label{rem:p0-pole-forward}
When a history leaves $B_\infty$, consecutive pencils $\mathcal R_k$ and
$\mathcal R_{k+1}$ are based at different points of $\partial_\infty\mathbb H^2$
and therefore have no common tangency point; the families cannot be nested, and
the pair carries a nonzero relative datum.  Quantifying that datum is not part
of this paper, whose scope fixes it to the poles and pencils themselves.  A
stabilizer-valued transport on the projective quotient that measures it is
constructed in Paper~IV; this remark is a prospective interface only and is not
used anywhere in the present paper.
\end{remark}
  at the end of
%   paper-0/sections/04-ripple-geometry.tex.
%
% Labels required from the host paper (all present):
%   eq:p0-mobius-map, def:p0-ripple-pencil, thm:p0-ripple-pencil-formula,
%   prop:p0-path-ripple-duality, eq:p0-projective-factorization,
%   eq:p0-bruhat, sec:p0-projective-unification.
% New labels introduced here (checked unique in the host paper):
%   subsec:p0-pole-covariance, lem:p0-pole-covariance,
%   prop:p0-pole-rigidity, rem:p0-pole-bruhat, rem:p0-pole-forward.
%
% Scope note.  Paper 0's scope file permits "the dual reading (poles,
% horocycle pullback, ripple pencils), reciprocals, infinity, fixed points,
% and the matrix synthesis", and excludes "projective bivaluations, quotient
% towers, condensation".  Everything below is therefore phrased in poles and
% in the Bruhat decomposition only; the quantitative comparison of two
% pencils tangent at different poles is left to Paper IV by a prospective
% reference.
%
% Claim status: Lemma and Proposition proved; the two Remarks record
% consequences and a prospective interface, not new theorems.
%
% Numbers when inserted in the current build of Paper 0 (Section 4 ends at
% statement 4.5): Lemma 4.6, Proposition 4.7, Remarks 4.8 and 4.9, and the
% subsection is 4.4, shifting the status subsection to 4.5.
% =====================================================================

\subsection{The pole of a prefix: covariance and rigidity}
\label{subsec:p0-pole-covariance}

Section~\ref{sec:p0-ripples} introduced the pole $p_F=F^{-1}(\infty)$ of a
single projective operator, and the ripple history of a word
$F_1,\dots,F_n$ carries the pole sequence $p_k=G_k^{-1}(\infty)$ of the prefixes
$G_k=F_k\circ\cdots\circ F_1$.  It was observed there that
\emph{the common tangency point of the pencil may move with $k$}.  The two
elementary statements below make that observation exact: they record how a
pole transforms under composition, and they identify the histories along which
it does not move.  Neither statement is new mathematics; they are recorded
because the mobility of the tangency point is otherwise left as a picture.

Throughout, a M\"obius map is as in \eqref{eq:p0-mobius-map}, and
$B_\infty=\Stab_{\PGL_2(K)}(\infty)$ is the affine Borel sector of
Section~\ref{sec:p0-projective-unification}.

\begin{lemma}[Covariance of the pole]
\label{lem:p0-pole-covariance}
Let $P,Q$ be nondegenerate M\"obius maps and put $p_P:=P^{-1}(\infty)$.  Then
\begin{equation}\label{eq:p0-pole-covariance}
  p_{P\circ Q}=Q^{-1}\bigl(p_P\bigr).
\end{equation}
Consequently
\begin{equation}\label{eq:p0-pole-constancy-criterion}
  p_{P\circ Q}=p_Q
  \quad\Longleftrightarrow\quad
  p_P=\infty
  \quad\Longleftrightarrow\quad
  P\in B_\infty .
\end{equation}
\end{lemma}

\begin{proof}
$(P\circ Q)^{-1}=Q^{-1}\circ P^{-1}$, so
$p_{P\circ Q}=(Q^{-1}\circ P^{-1})(\infty)=Q^{-1}\bigl(p_P\bigr)$, which is
\eqref{eq:p0-pole-covariance}.  For the criterion: since $Q^{-1}$ is injective
on the projective line, $Q^{-1}(p_P)=Q^{-1}(\infty)$ holds exactly when
$p_P=\infty$, and $P^{-1}(\infty)=\infty$ holds exactly when
$P(\infty)=\infty$, that is $C_P=0$, that is $P\in B_\infty$.
\end{proof}

Lemma~\ref{lem:p0-pole-covariance} says that the pole moves by the action of the
\emph{next} operator on the projective line: the new pole is the old one moved by
$Q^{-1}$.  It also says that a step leaves the pole unchanged exactly when the
step is affine.  Applied to a history, the two statements give the rigidity
criterion.

\begin{proposition}[Rigidity of the pole along a history]
\label{prop:p0-pole-rigidity}
Let $F_1,\dots,F_n$ be nondegenerate M\"obius sections in chronological order,
$G_k=F_k\circ\cdots\circ F_1$, and $p_k=G_k^{-1}(\infty)$.  The following are
equivalent:
\begin{enumerate}
  \item $p_k=p_{k-1}$ for $k=1,\dots,n$;
  \item $F_k\in B_\infty$ for $k=1,\dots,n$, i.e. every lower-left entry
  vanishes;
  \item $p_k=\infty$ for $k=1,\dots,n$.
\end{enumerate}
Equivalently, the pole sequence is not constantly $\infty$ exactly when the
history leaves the affine sector, and then the first index at which $C_k\ne0$ is
also the first index at which the tangency point is finite.
\end{proposition}

\begin{proof}
Apply Lemma~\ref{lem:p0-pole-covariance} with $P=F_k$ and $Q=G_{k-1}$: the step
$G_{k-1}\to G_k$ leaves the pole unchanged exactly when $F_k\in B_\infty$.  This
gives (1)$\Leftrightarrow$(2).  Since $G_0=\id$ and $p_0=G_0^{-1}(\infty)
=\infty$, the equivalence (2)$\Leftrightarrow$(3) follows by induction: if
every $F_k$ is affine then every prefix is affine, so $C_k=0$ and
$p_k=\infty$; conversely if some $C_k\ne0$ then $p_k=-D_k/C_k\ne\infty$ by the definition of the pole in
Section~\ref{sec:p0-ripples}.
\end{proof}

\begin{remark}[The pole is read off one factor of the factorization]
\label{rem:p0-pole-bruhat}
When $C_k\ne0$, the factorization \eqref{eq:p0-projective-factorization} of
$G_k$ reads
\[
  G_k
  =T_{A_k/C_k}\circ D_{\Delta_k/C_k^{2}}\circ J\circ T_{D_k/C_k},
\]
and the pole is
$p_k=G_k^{-1}(\infty)=-D_k/C_k$, i.e. the negative of the translation parameter
of the innermost factor.  Thus of the four factors, the one that determines the
tangency point of the pencil is the rightmost, while the value orbit
$G_k(z_0)$ depends on the whole word.  This is the precise sense in which the
pole and the point are independent data of one prefix, and it is the reading of
the cell decomposition \eqref{eq:p0-bruhat} in terms of tangency: the pole is
$\infty$ exactly in the affine cell, and it is finite exactly in the projective
cell.
\end{remark}

\begin{remark}[Prospective interface]
\label{rem:p0-pole-forward}
When a history leaves $B_\infty$, consecutive pencils $\mathcal R_k$ and
$\mathcal R_{k+1}$ are based at different points of $\partial_\infty\mathbb H^2$
and therefore have no common tangency point; the families cannot be nested, and
the pair carries a nonzero relative datum.  Quantifying that datum is not part
of this paper, whose scope fixes it to the poles and pencils themselves.  A
stabilizer-valued transport on the projective quotient that measures it is
constructed in Paper~IV; this remark is a prospective interface only and is not
used anywhere in the present paper.
\end{remark}
  at the end of
%   paper-0/sections/04-ripple-geometry.tex.
%
% Labels required from the host paper (all present):
%   eq:p0-mobius-map, def:p0-ripple-pencil, thm:p0-ripple-pencil-formula,
%   prop:p0-path-ripple-duality, eq:p0-projective-factorization,
%   eq:p0-bruhat, sec:p0-projective-unification.
% New labels introduced here (checked unique in the host paper):
%   subsec:p0-pole-covariance, lem:p0-pole-covariance,
%   prop:p0-pole-rigidity, rem:p0-pole-bruhat, rem:p0-pole-forward.
%
% Scope note.  Paper 0's scope file permits "the dual reading (poles,
% horocycle pullback, ripple pencils), reciprocals, infinity, fixed points,
% and the matrix synthesis", and excludes "projective bivaluations, quotient
% towers, condensation".  Everything below is therefore phrased in poles and
% in the Bruhat decomposition only; the quantitative comparison of two
% pencils tangent at different poles is left to Paper IV by a prospective
% reference.
%
% Claim status: Lemma and Proposition proved; the two Remarks record
% consequences and a prospective interface, not new theorems.
%
% Numbers when inserted in the current build of Paper 0 (Section 4 ends at
% statement 4.5): Lemma 4.6, Proposition 4.7, Remarks 4.8 and 4.9, and the
% subsection is 4.4, shifting the status subsection to 4.5.
% =====================================================================

\subsection{The pole of a prefix: covariance and rigidity}
\label{subsec:p0-pole-covariance}

Section~\ref{sec:p0-ripples} introduced the pole $p_F=F^{-1}(\infty)$ of a
single projective operator, and the ripple history of a word
$F_1,\dots,F_n$ carries the pole sequence $p_k=G_k^{-1}(\infty)$ of the prefixes
$G_k=F_k\circ\cdots\circ F_1$.  It was observed there that
\emph{the common tangency point of the pencil may move with $k$}.  The two
elementary statements below make that observation exact: they record how a
pole transforms under composition, and they identify the histories along which
it does not move.  Neither statement is new mathematics; they are recorded
because the mobility of the tangency point is otherwise left as a picture.

Throughout, a M\"obius map is as in \eqref{eq:p0-mobius-map}, and
$B_\infty=\Stab_{\PGL_2(K)}(\infty)$ is the affine Borel sector of
Section~\ref{sec:p0-projective-unification}.

\begin{lemma}[Covariance of the pole]
\label{lem:p0-pole-covariance}
Let $P,Q$ be nondegenerate M\"obius maps and put $p_P:=P^{-1}(\infty)$.  Then
\begin{equation}\label{eq:p0-pole-covariance}
  p_{P\circ Q}=Q^{-1}\bigl(p_P\bigr).
\end{equation}
Consequently
\begin{equation}\label{eq:p0-pole-constancy-criterion}
  p_{P\circ Q}=p_Q
  \quad\Longleftrightarrow\quad
  p_P=\infty
  \quad\Longleftrightarrow\quad
  P\in B_\infty .
\end{equation}
\end{lemma}

\begin{proof}
$(P\circ Q)^{-1}=Q^{-1}\circ P^{-1}$, so
$p_{P\circ Q}=(Q^{-1}\circ P^{-1})(\infty)=Q^{-1}\bigl(p_P\bigr)$, which is
\eqref{eq:p0-pole-covariance}.  For the criterion: since $Q^{-1}$ is injective
on the projective line, $Q^{-1}(p_P)=Q^{-1}(\infty)$ holds exactly when
$p_P=\infty$, and $P^{-1}(\infty)=\infty$ holds exactly when
$P(\infty)=\infty$, that is $C_P=0$, that is $P\in B_\infty$.
\end{proof}

Lemma~\ref{lem:p0-pole-covariance} says that the pole moves by the action of the
\emph{next} operator on the projective line: the new pole is the old one moved by
$Q^{-1}$.  It also says that a step leaves the pole unchanged exactly when the
step is affine.  Applied to a history, the two statements give the rigidity
criterion.

\begin{proposition}[Rigidity of the pole along a history]
\label{prop:p0-pole-rigidity}
Let $F_1,\dots,F_n$ be nondegenerate M\"obius sections in chronological order,
$G_k=F_k\circ\cdots\circ F_1$, and $p_k=G_k^{-1}(\infty)$.  The following are
equivalent:
\begin{enumerate}
  \item $p_k=p_{k-1}$ for $k=1,\dots,n$;
  \item $F_k\in B_\infty$ for $k=1,\dots,n$, i.e. every lower-left entry
  vanishes;
  \item $p_k=\infty$ for $k=1,\dots,n$.
\end{enumerate}
Equivalently, the pole sequence is not constantly $\infty$ exactly when the
history leaves the affine sector, and then the first index at which $C_k\ne0$ is
also the first index at which the tangency point is finite.
\end{proposition}

\begin{proof}
Apply Lemma~\ref{lem:p0-pole-covariance} with $P=F_k$ and $Q=G_{k-1}$: the step
$G_{k-1}\to G_k$ leaves the pole unchanged exactly when $F_k\in B_\infty$.  This
gives (1)$\Leftrightarrow$(2).  Since $G_0=\id$ and $p_0=G_0^{-1}(\infty)
=\infty$, the equivalence (2)$\Leftrightarrow$(3) follows by induction: if
every $F_k$ is affine then every prefix is affine, so $C_k=0$ and
$p_k=\infty$; conversely if some $C_k\ne0$ then $p_k=-D_k/C_k\ne\infty$ by the definition of the pole in
Section~\ref{sec:p0-ripples}.
\end{proof}

\begin{remark}[The pole is read off one factor of the factorization]
\label{rem:p0-pole-bruhat}
When $C_k\ne0$, the factorization \eqref{eq:p0-projective-factorization} of
$G_k$ reads
\[
  G_k
  =T_{A_k/C_k}\circ D_{\Delta_k/C_k^{2}}\circ J\circ T_{D_k/C_k},
\]
and the pole is
$p_k=G_k^{-1}(\infty)=-D_k/C_k$, i.e. the negative of the translation parameter
of the innermost factor.  Thus of the four factors, the one that determines the
tangency point of the pencil is the rightmost, while the value orbit
$G_k(z_0)$ depends on the whole word.  This is the precise sense in which the
pole and the point are independent data of one prefix, and it is the reading of
the cell decomposition \eqref{eq:p0-bruhat} in terms of tangency: the pole is
$\infty$ exactly in the affine cell, and it is finite exactly in the projective
cell.
\end{remark}

\begin{remark}[Prospective interface]
\label{rem:p0-pole-forward}
When a history leaves $B_\infty$, consecutive pencils $\mathcal R_k$ and
$\mathcal R_{k+1}$ are based at different points of $\partial_\infty\mathbb H^2$
and therefore have no common tangency point; the families cannot be nested, and
the pair carries a nonzero relative datum.  Quantifying that datum is not part
of this paper, whose scope fixes it to the poles and pencils themselves.  A
stabilizer-valued transport on the projective quotient that measures it is
constructed in Paper~IV; this remark is a prospective interface only and is not
used anywhere in the present paper.
\end{remark}
